\section{Multi-Hop Reasoning and Query Planning}
\label{sec:gr-multihop}

Many realistic questions require combining several facts that no single passage
contains, a setting formalized by multi-hop question answering benchmarks that
demand reasoning over multiple supporting documents~\cite{yang2018hotpotqa}. A
knowledge graph is naturally suited to this task because a multi-hop question maps
onto a path or subgraph, and following labelled edges makes each reasoning step
explicit and traceable~\cite{peng2024graphrag}. This explicitness is a key
advantage of graph-structured retrieval over flat vector search, whose similarity
scores give no principled way to chain intermediate facts
together~\cite{peng2024graphrag}.

Because such questions must be decomposed before they can be traversed, Graph RAG
systems increasingly plan retrieval: complex queries are broken into simpler
sub-queries whose results are combined, and language models are used to drive this
planning over the graph~\cite{peng2024graphrag}. Coupling the reasoning ability of
language models with the structured, verifiable knowledge of a graph is exactly
the complementarity emphasized in roadmaps for unifying the two
paradigms~\cite{pan2024unifying}.
